% !TeX root = ../main.tex

\section{Grundlagen}
\subsection{Smart Grids: physische Ebene sowie Informations- und Kommunikationstechnik}
Moderne Energiesysteme befinden sich in einem tiefgreifenden Wandel, der durch den verstärkten Einsatz erneuerbarer Energien, dezentrale Erzeugungsstrukturen und einer zunehmenden Volatilität im Netzbetrieb gekennzeichnet ist. Um diesen Herausforderungen erfolgreich begegnen zu können, entwickeln sich traditionelle Stromnetze zu sogenannten Smart Grids weiter. Diese kombinieren reale Infrastruktur mit digitalen Technologien für Überwachung, Steuerung und Effizienz (Gungor et al., 2013). Die physische Ebene besteht dabei aus Erzeugungsanlagen, Umspannwerken, Leitungen, Transformatoren und Lasten, die gemeinsam die zusammen das Grundgerüst der Netzstruktur bilden. Ergänzt wird sie durch eine digitale Ebene, welche aus Messsystemen, Kommunikationsnetzen, Regelalgorithmen und Datenplattformen besteht. Letztere fördert Austauschprozesse zwischen Anbietern und Nutzern (Fang et al., 2012).
Eine Schlüsselkomponente dieser digitalen Infrastruktur ist die sog. Advanced Metering Infrastructure (AMI), welche mithilfe intelligenter Zähler den Energieverbrauch in Echtzeit erfasst und damit zur Netzzustandsbewertung beiträgt. Weitere wichtige Komponenten sind SCADA-Systeme, Sensorik zur Überwachung von Spannungen und Frequenzen sowie Kommunikationsprotokolle wie IEC 61850, die den Datenaustausch im Netz standardisieren (NIST, 2014). Mit diesen Bausteinen lässt sich das Netz genauer beobachten, Lasten effizienter steuern, Fehler rascher identifizieren. Allerdings führt diese Vernetzung zwangsläufig dazu, dass digitale und reale Systeme enger verflochten sind, was wiederum neue Schwachstellen begünstigt. Angriffe auf Netzwerke oder Fehler in Programmen können direkte Auswirkungen auf technische Systeme haben und Störungen entstehen, die sich schrittweise verstärken lassen (Buldyrev et al., 2010). Bei der Untersuchung der Widerstandsfähigkeit ergibt sich daher ein zentrales Problem: physikalische wie elektronische Vernetzungen müssen einzeln erfasst werden; nur zusammen erklären sie das Gesamtsystemverhalten.
\subsection{Komplexe Netzwerke: Knoten, Kanten, typische Topologien}
Komplexe Netze zeigen vereinfacht, wie echte Systeme funktionieren – dabei werden Bausteine mit ihren Verbindungen durch Grafiken dargestellt. Eine solche Struktur enthält Punkte, sogenannte Knoten, verbunden über Linien, genannt Kanten, welche Austauschprozesse abbilden (Newman, 2018). In intelligenten Stromnetzen stehen diese Knoten meist für technische Komponenten wie Umspannwerke, Transformatorstationen, Steuerzentralen oder lokale Energiequellen. Die verbindenden Elemente hingegen symbolisieren entweder reale Leitungen. Entscheidend dafür, wie gut das Gesamtsystem Störungen aushält, ist nicht nur die Menge der beteiligten Einheiten, sondern auch deren Vernetzungsart.
\subsubsection{Knoten}
Ein Knoten bildet die Basis eines Netzwerks. Abhängig vom Kontext handelt es sich um ein Gerät, eine Stelle in einer Struktur oder einen gedanklichen Punkt. Die Zahl der Linien, die von ihm wegführen, heißt Grad des Knotens. Wichtig ist dieser Wert, weil er zeigt, wie bedeutend ein Knoten ist. Besonders viele Verbindungen haben oft wichtige Knoten im Stromnetz, das macht sie empfindlich.
\subsubsection{Kanten} 
Kanten zeigen, wie Knoten miteinander verbunden sind. In technischen Systemen, etwa bei Smart Grids, dienen sie dem Transport von Energie oder Daten. Zu ihren zentralen Merkmalen gehören sogenannte gewichtete Kanten. Dieses Gewicht steht im Stromnetz beispielsweise für Parameter wie Leitungsbelastbarkeit oder Widerstand. Darüber hinaus spielt die Richtung (gerichtet/ungerichtet) eine Rolle. Bei rein struktureller Betrachtung gelten Stromnetze als ungerichtet; sobald Leistungsflüsse ins Gewicht fallen, werden sie dagegen als gerichtet angesehen.
Knoten mit Kanten prägen die Struktur eines Netzwerks, was sein Verhalten bei Belastung beeinflusst, ebenso wie die Reaktion auf Störungen.
\subsubsection{Zentralitätsmaße}
Zentralitätsmaße quantifizieren die Bedeutung eines Knotens. Beispiele dafür sind die folgenden drei Maße.
	Betweenness Centrality: Sie misst, wie häufig ein Knoten auf den kürzesten Pfaden zwischen anderen Knoten liegt. Das ist wichtig für die Flusskontrolle. Sie wird durch folgenden Ausdruck gemessen:
\[
g(v) = \sum_{\substack{s \neq v \neq t}} \frac{\sigma_{st}(v)}{\sigma_{st}}
\]

σ_st ist dabei die die Gesamtzahl der kürzesten Wege vom Knoten s zum Knoten t und σ_st die Anzahl derjenigen dieser Wege, die durch den Knoten v verlaufen (wobei v kein Endpunkt des Weges ist).
	Closeness Centrality: Sie misst, wie nah ein Knoten durchschnittlich zu allen anderen liegt. Das erfolgt folgendermaßen: 
\[
C(x) = \frac{N - 1}{\sum_{y} d(y, x)}
\]

N ist in der Abbildung dabei die Anzahl der Knoten im Graphen. Die Normalisierung der Closeness (Nähe-Zentralität) erleichtert den Vergleich von Knoten in Graphen unterschiedlicher Größe. Bei sehr großen Graphen wird das „minus eins“ in der Normalisierung vernachlässigbar und daher häufig weggelassen.
	Degree Centrality: Die Degree Centrality bezeichnet die Bedeutung eines Knotens in einem Netzwerk anhand der Anzahl seiner direkten Verbindungen. Sie entspricht dem Knotengrad und misst, wie viele Kanten von einem Knoten ausgehen. Ein Knoten mit hoher Degree Centrality hat viele Nachbarn und spielt daher eine wichtige Rolle für die lokale Konnektivität des Netzwerks. In ungerichteten Netzwerken entspricht die Degree Centrality der Anzahl aller benachbarten Knoten. (Freeman, 1979).
\subsubsection{Topologien}
Komplexe Netzwerke können auf unterschiedliche Weise strukturiert sein. Drei grundlegende Modelle haben sich bisher zur Beschreibung und Simulation realer Systeme etabliert: Erdős–Rényi-Netzwerke, Watts–Strogatz-Small-World-Netzwerke und Barabási–Albert-skalenfreie Netzwerke. Diese Topologien unterscheiden sich hinsichtlich Verbindungsstruktur, Gradverteilung, Clusterkoeffizient und typischer Pfadlängen und damit auch hinsichtlich ihrer Resilienz.
Das Erdős–Rényi-Netzwerk, ER-Modell, (Erdős & Rényi, 1959) erzeugt ein Netzwerk, indem zwischen allen möglichen Knotenpaaren jeweils mit einer festen Wahrscheinlichkeit p eine Kante gezogen wird. Diese zufällige Verbindungsstrategie hat mehrere Konsequenzen:
	Die Gradverteilung ist binomial- bzw. in großen Netzen näherungsweise Poisson-verteilt und die meisten Knoten haben ähnliche Gradwerte.
	Der Clusterkoeffizient ist vergleichsweise gering, da nur wenige Dreiecksbeziehungen entstehen.
	Die Pfadlänge ist typischerweise klein (logarithmisch mit Netzwerkgröße).
	Die Resilienz ist moderat robust gegenüber zufälligen Ausfällen, aber vulnerabel gegenüber gezielten Angriffen, da keine systematische Struktur existiert, die kritische Knoten schützt.
ER-Netzwerke bilden eine wichtige Baseline, aber reale Stromnetze weisen häufig strukturierte Muster auf, die durch dieses Modell nicht vollständig beschrieben werden können.
Das Small-World-Netzwerk-Modell (Watts & Strogatz, 1998) kombiniert Merkmale regulärer und zufälliger Netzwerke. Die Konstruktion erfolgt durch Ausgang eines regulären Rings, in dem Knoten mit ihren nächsten Nachbarn verbunden sind, und anschließend eine zufällige „Umverteilung“ einzelner Kanten passiert.
Charakteristische Eigenschaften:
	Eine hohe lokale Clusterbildung - also eine lokale Redundanz - wird geschaffen. Diese Knoten haben viele gemeinsame Nachbarn.
	Die durchschnittliche Pfadlänge ist eher kurz. Dadurch existiert eine effiziente Erreichbarkeit im gesamten Netzwerk.
	Die Resilienz ist durch hohe lokale Clusterung vergleichsweise hoch und damit das Netzwerk widerstandsfähiger gegenüber Ausfällen. Jedoch ist es anfällig für Störungen, wenn die wenigen „Shortcut“-Kanten ausfallen, die die Pfadlänge stark reduzieren.
Small-World-Strukturen treten in vielen Energie- und Informationsnetzen auf, da sie Effizienz mit lokaler Stabilität verbinden.
Von besonderer Bedeutung sind skalenfreie Netzwerke nach Barabási und Albert, deren Knotengradverteilung einem Potenzgesetz folgt. Das BA-Modell (Albert, Jeong & Barabási, 2000) basiert auf „präferentieller Anbindung“ von Knoten. Es werden nacheinander Knoten hinzugefügt und sie werden bevorzugt mit bereits gut vernetzten Knoten verbunden („the rich get richer“). Dadurch entstehen sog. Hubs.
Eigenschaften:
	Die Gradverteilung folgt einem Potenzgesetz (scale-free). Dadurch entstehen wenige Hubs mit vielen Knoten mit geringem Grad.
	Der Clusterkoeffizient ist vergleichsweise gering, kann aber je nach Erweiterung variieren.
	Die Pfadlängen ist sehr kurz, da Hubs viele Wege verbinden (ultra-small world).
	Die Resilienz ist hoch, da das Netzwerk sehr robust gegen zufällige Ausfälle ist (da diese meist kleine Knoten betreffen). Es ist aber stark anfällig gegenüber gezielten Angriffen auf Hubs, da das Netzwerk leicht fragmentieren kann.
Viele reale Infrastrukturnetze, einschließlich Kommunikationssysteme, Internet-Backbones und Teile moderner Stromnetze, zeigen skalenfreie Eigenschaften. Daher ist das BA-Modell besonders relevant für Analysen der Verwundbarkeit kritischer Infrastrukturen.

\subsection{Zentrale Metriken} 
Im Folgenden werden die Metriken, welche zentral für die Analyse von Netzwerkstrukturen sind, erläutert und dargelegt, welche Relevanz diese jeweils für Smart-Grids haben. 
\subsubsection{Gradverteilung P(k)}
Die Gradverteilung P(k) ist die Wahrscheinlichkeit, dass ein Knoten eines Netzwerkes Grad k hat und liegt zwischen 0 und 1. Die Gradverteilung eines Netzwerks zeigt die Beschaffenheit deren Gesamtstruktur auf und kann somit der Feststellung damit einhergehender, typischer Eigenschaften dienen. Für Smart Grids ist die Feststellung solcher typischen Eigenschaften insofern relevant, als dass hierdurch kritische Hubs identifiziert werden können, wie z.B. große Umspannwerke mit einer besonders hohen Anzahl an Leitungsverbindungen. Die Gradverteilung von Stromnetzwerken weist die grundsätzlichen Bedingungen der Skalenfreiheit auf: Zum Einen ist durch stetigen Ausbau des Stromnetzes das Wachstum als Voraussetzung gegeben. Ebenso werden die Kraftwerke vorzugsweise über stärker vernetzte Umspannwerke angebunden, anstatt perifer dezentral geleitet zu werden. Somit entsteht die typische Gradverteilung, die dem Potenzgesetz entspricht P(k) ~ k-γ. Bei dieser Verteilung wird die Wahrscheinlichkeit P(k) bei wachsendem k immer kleiner. Im Falle von Stromnetzwerken, bspw. in der westlichen USA, liegt γ ca. bei 4 und somit an der Grenze der Eigenschaften der Skalenfreiheit : die Anzahl großer Hubs ist vergleichsweise gering und stößt mit einer Gesamtzahl von einigen tausend Kraftwerke relativ schnell an eine natürliche Obergrenze. 
\subsubsection{Clustering-Koeffizient C} 
Der Clustering-Koeffizient C beschreibt, ob zwei Nachbarn eines Knotens ebenfalls füreinander Nachbarn sind  und damit die Ausprägung der Transitivität des Netzwerkes : 
\[
C = \frac{\text{Anzahl geschlossener Pfade der L\"ange 2}}{\text{Anzahl Pfade der L\"ange 2}}
\]
Bei einer sehr großen Anzahl an Knoten N eines Netzwerks wird der Clustering-Koeffizient sehr klein. Diese Ausprägung eines Extremfalls ist bei Stromnetzen aufgrund der begrenzten Knotenanzahl jedoch nicht gegeben. Relevant wird der Clustering-Koeffizient dann, wenn einzelne Knoten und Verbindungen ausfallen. Der Koeffizient gibt an, wie groß die Redundanz im Netzwerk ist. Ist ein hoher Clustering-Koeffizient gegeben, so stehen im Netzwerk mehrere alternative Pfade bei Ausfällen zur Verfügung, was für Robustheit des Netzwerks gegen lokale Ausfälle sorgt. 
\subsubsection{Mittlere Pfadlänge l, Durchmesser und Effizienz}
Die mittlere Pfadlänge beschreibt die mittlere kürzeste Entfernung zwischen Knoten. Dieser Kennwert ist relevant, um die Ausbreitung von Lastfluss oder Störungen zu beschreiben und um die Wiederherstellungszeit einzuschätzen . Die mittlere Pfadlänge lässt sich wiefolgt berechen:
\[
\ell = \frac{1}{N(N-1)} \sum_{i j} d_{i j}
\]
Wobei dij die durchschnittliche kürzeste Distanz zwischen den Knoten i und j eines Netzwerks bezeichnet . 
\subsection{Robustheit, Resilienz, Fehlertoleranz}
In der Literatur zu komplexen Netzwerken werden die Begriffe Robustheit, Fehlertoleranz und Resilienz häufig gemeinsam verwendet, obwohl diese verschiedene Schwerpunkte fokussieren . Im den folgenden Unterkapiteln wird eine Abgrenzung im Kontext von Stromnetzwerken vorgenommen, basierend auf der Übersichtsarbeit von Xiaoyu Qi und Gang Mei (2024). 
\subsubsection{Robustheit}
Robustheit beschreibt die Fähigkeit eines Netzwerks, seine Strukturen und Funktionen trotz Störungen weitgehend gleichbleibend aufrechtzuerhalten . Robustheit kann als Teilaspekt von Resilienz betrachtet werden  und bezieht sich hauptsächlich darauf, wie widerstandsfähig ein Netzwerk in seiner Struktur selbst gegenüber Störungen ist. Im Bereich von Stromnetzwerken, die zu den technologischen Netzwerken zählen, wird Robustheit als die Resistenz gegenüber Ausfällen bezeichnet. Der Kernfrage der Robustheit besteht darin, wie gut die Netzwerkfunktion unmittelbar nach einer Störung erhalten bleibt. In diesem Sinne gilt ein Stromnetz als robust, sofern der Ausfall eines einzelnen Umspannwerks keine weiterführenden Stromausfälle verursacht. 
Gemessen wird die Robustheit typischer Weise über die größte zusammenhängende Komponente (GCC). Hierbei werden nach und nach Knoten oder Kanten entweder zufällig oder gezielt aus dem Netzwerk entfernt und nach jedem Schritt die Größe des größten funktionierenden Netzteils gemessen : 
\[
\mathrm{Rob} = \frac{1}{N} \sum_{q=1}^{N} S(q)
\]
Wobei Rob die Gesamtrobustheit eines Netzwerks bezeichnet und S(q) die relative Größe der größten Komponente nach q entfernten Knoten darstellt . 
\subsubsection{Fehlertoleranz}
Die Fehlertoleranz bezeichnet die Fähigkeit eines Systems, trotz Fehler einzelner Komponenten weiterhin funktionsfähig zu bleiben und ohne die Fehler zwingend reparieren zu müssen. Während Stromnetze im Bereich der technologischen Netze angesiedelt sind, bezieht sich die Fehlertoleranz mehr auf Informationsnetzwerke , zu denen Kommunikations- und Internetnetzwerke zählen . Dabei wird stärker auf Hard- und Software Bezug genommen und die Fähigkeit beschrieben, Fehler strukturell oder algorithmisch zu kompensieren, z.B. durch Redundanzen oder Umleitung von Flüssen. Die Kernfrage der Fehlertoleranz ist, wie gut das System während eines auftretenden Fehlers weiterarbeiten kann. So gilt ein Stromnetz als fehlertolerant, wenn bei einem Leitungsausfall automatisch alternative Wege für den Stromfluss gewählt werden. 
\subsubsection{Resilienz}
Die Resilienz umfasst die Widerstandsfähigkeit eines Netzwerks im Störungsfall, die Fähigkeit zur Begrenzung der Auswirkungen und die Fähigkeit zur Wiederaufnahme eines akzeptablen Funktionsniveaus. Dementsprechend ist die Resilienz ein übergreifendes Konzept, welches die Robustheit, Fehlertoleranz, Anpassungs- und Erholungsfähigkeit beinhaltet. Ein Smart Grid ist beispielsweise dann resilient, wenn es nach einem extremen Wetterereignis schnell wieder auf ein normales Versorgungsniveau zurückkehren kann . 
Gemessen werden kann die Resilienz anhand der Fläche unter der Leistungs-Zeit-Kurve: 
\[
\mathrm{Res} = \int_{t_0}^{t_1} Q(t)\,\mathrm{d}t
\]
Wobei Q(t) die Systemleistung relativ zur Norm ist und Res beschreibt, wie schnell und wie gut die Erholung erfolgt. Die Leistung kann die Konnektivität (GCC), Flusseffizienz, Versorgungsgrad oder Kommunikationsfähigkeit sein und kann berechnet werden durch 
\[
Q(t) = \frac{P(t)}{TP(t)}
\], wobei P(t) die gemessene und TP(t) die Zielleistung ist, die ohne Störung zu erwarten wäre. 
Bei mehreren Messdurchgängen wird der Erwartungswert von 
\[
\frac{P(t)}{TP(t)}
\] (also der Mittelwert der Messdurchgänge) berechnet. Die Kernfrage der Resilienz ist, wie viel der ursprünglichen Leistung noch vorhanden ist. 
